\section{Concluzii}

Lucrarea a acoperit doua directii complementare: implementarea unui predictor de
ramificatii \textbf{OGEHL} (Optimized Geometric History Length) ca librarie C
portabila si integrarea sa in doua medii de executie distincte --- o aplicatie
grafica independenta si simulatorul \textbf{SimpleScalar sim-bpred} --- urmate de
o evaluare cantitativa pe patru benchmark-uri din suita \textbf{SPEC CPU95}.

\begin{itemize}
    \item \textbf{Corectitudinea implementarii.}
          Rezultatele pe benchmark-ul \textbf{mgrid} (99.19\%--99.21\% acuratete
          pe toate cele 9 configuratii testate) confirma ca implementarea OGEHL
          este corecta si stabila: un program cu flux de control regulat si
          predictibil este ghicit corect in proportie covarsitoare, indiferent
          de parametri. Aceasta uniformitate serveste ca test de sanity pentru
          intreaga infrastructura de simulare.

    \item \textbf{Sensibilitatea la parametri depinde de natura aplicatiei.}
          Pe \textbf{applu} (dinamica fluidelor, floating-point intensiv),
          marimea tabelelor este factorul dominant: configuratii mici
          ($S=512$) produc acurateti de 68--75\% din cauza aliasing-ului masiv,
          in timp ce $S=2048$ ridica acuratetea la $\approx$98\%. Pe
          \textbf{compress} (compresie LZW, ramificatii dependente de date),
          lungimea contoarelor devine critica: trecerea de la $k=4$ la $k=5$
          biti produce un salt de $\approx$6\%, indiferent de $N$ si $S$. Pe
          \textbf{ijpeg} (compresie JPEG), variatia este moderata (0.74\%),
          cu marimea tabelelor ca factor principal si lungimea contoarelor fara
          efect masurabil.

    \item \textbf{Configuratia optima general-valabila este c5.}
          Configuratia $N=7$, $S=2048$, $k=5$ atinge acuratetea maxima pe
          applu (98.37\%), este near-optimal pe mgrid si ijpeg si produce al
          doilea cel mai bun rezultat pe compress (89.37\%, la 0.53\% de
          optim), cu un consum de memorie de $\approx$90~KB --- de 4 ori mai
          eficienta decat configuratiile cu $S=8192$.

    \item \textbf{Numarul de tabele are impact conditionat.}
          Cresterea de la $N=7$ la $N=12$ fara cresterea proportionala a $S$
          \textit{scade} acuratetea pe applu (c7: 95.58\% vs.\ c4: 97.36\%),
          confirmand ca fragmentarea bugetului de memorie in mai multe tabele
          mai mici amplifica aliasing-ul per tabel si anuleaza castigul din
          contextele de istorie mai lungi.

    \item \textbf{Mecanismele adaptive ale OGEHL isi dovedesc utilitatea.}
          Ajustarea dinamica a pragului $\theta$ si a lungimii istoriei
          (Dynamic History Length Fitting) permit predictorului sa se
          adapteze la profilul de executie fara retuning manual, ceea ce
          explica acuratetea ridicata pe benchmark-uri cu caractere diferite
          folosind aceiasi parametri.

    \item \textbf{SimpleScalar ramane o platforma viabila pentru cercetare.}
          Integrarea OGEHL in \textit{sim-bpred} a necesitat doar adaugarea
          unui nou tip de predictor in \texttt{bpred.c} si \texttt{bpred.h},
          fara modificari in pipeline-ul simulatorului. Aceasta confirma
          modularitatea arhitecturii SimpleScalar si accesibilitatea sa pentru
          extinderi academice.
\end{itemize}

In ansamblu, lucrarea demonstreaza ca OGEHL reprezinta o solutie eficienta si
configurabila pentru predictia de ramificatii, cu un raport acuratete/memorie
superior predictorilor clasici din SimpleScalar (bimodal, 2-level). Directii
viitoare de cercetare includ evaluarea pe suite mai largi (SPEC CPU2000/2006),
analiza consumului de putere in functie de configuratii.
